\section{The Eucharistic Foretaste}

Christ gives the pilgrim Church a pledge of the life she awaits. In the Eucharist he gives himself, and the communion promised in glory is already begun sacramentally. The Catechism calls the Eucharist a pledge of future glory and directs its meaning toward the Lord's return (CCC 1402--1405). The same Savior who will be our eternal joy comes to us under the sacramental signs. His presence is real while its manner remains veiled. Faith receives what sight does not yet disclose, and the heart is nourished for a fulfillment it cannot attain by its own strength.

In receiving the Eucharist, we receive the same Christ whom we hope to see in glory. His death and Resurrection have opened heaven, and the Eucharistic sacrifice makes present his saving self-offering. The sacrament therefore joins remembrance and expectation. The Church remembers the Paschal mystery while awaiting its consummation in the kingdom. She awaits the fulfillment of what the Lord has accomplished, while the risen Christ continues to give the fruits of his sacrifice to his members.

Irenaeus reasons from the Eucharist to the resurrection of the body. Created bread and wine belong within the sacrament by which Christ nourishes his people; the bodies that receive him are destined for resurrection (\emph{Against Heresies} V.2.2--3). Our dependence upon material food and our dependence upon divine life meet within God's saving generosity. We receive this gift as embodied members of Christ, in the flesh whose salvation Irenaeus defends. The Creator's handiwork is taken into the service of a redemption that will reach its bodily fullness.

The communion is ecclesial as well. \emph{Lumen gentium} 50 places our closest sacramental union with the heavenly Church in the Eucharistic sacrifice. The saints whom we commemorate are joined with us in the praise of the triune God. Mary, Joseph, the apostles, martyrs, and all the blessed belong to the communion in which the pilgrim Church worships. The Mass can be celebrated in a poor or nearly empty church, amid distraction and human weakness, and still possess this heavenly breadth. What faith knows of its communion far exceeds what the eye can count in the room.

Eucharistic communion also makes a claim upon our conduct toward one another. We receive the Lord whose will is the unity of his people in charity. It is possible to desire spiritual consolation while remaining unwilling to bear the inconvenience of a fellow member. The sacrament calls that divided desire toward healing. A person who needs our forgiveness, our material help, or simply our patient attention cannot be excluded from the meaning of communion. We can bring our reluctance to Christ and seek the help we need to fulfill that inconvenient duty.

The foretaste does not always feel like one. A distracted reception, a long period of dryness, or the grief of a difficult season may leave little sensible joy. The sacrament's truth is not established by our experience of consolation. The word \emph{pledge} directs attention to the fidelity of the giver. Christ commits himself to his people, nourishes their life of grace, and draws them toward their end. We can give thanks for his presence even when thanksgiving is quiet and laborious. Such faith honors the difference between the gift itself and the feelings by which we sometimes recognize its sweetness.

Adoration teaches the same patience. To remain before Christ is to consent to a presence whose depth we do not command. We bring our needs, but we also learn to be with him for his own sake. The familiar pressure to make every interval productive can loosen in that act. A creature's highest good is communion with God; time spent attending to him is already ordered toward its end. This does not relieve us of duties outside the church. It restores their order by joining them to the Lord from whom they receive their meaning and in whom their charity will endure.

At the threshold of death the Eucharistic promise takes on the intimacy of food for the journey. The Christ whom the Christian has received throughout life remains the one toward whom he goes. Even when illness has made bodily reception impossible, desire for him can remain within the Church's prayer and care. The final hope rests in his mercy and fidelity. The sacramental life has been preparing the person to receive him wholly: the God hidden under the signs is the God who will make himself known face to face, and the communion begun in faith is ordered to an unending life of love.

